\section{Empresas de datos, Meta/Scale y organización del talento}

Lo anterior trató de la tecnología de datos para robots; esta sección se ocupa de por qué la capacidad de producir datos se ha vuelto una capacidad estratégica de las grandes empresas. En la entrevista, Zhang Xiaojun (张小珺) y Xie Chen (谢晨) comentan la adquisición de Scale AI por parte de Meta a un precio elevado y el papel de Alexandr Wang. Lo importante aquí no es el chisme de la operación en sí, sino una conclusión más amplia: en la siguiente etapa, la competencia en IA no dependerá solo del poder de cómputo y de los modelos, sino también de la capacidad de organizar la producción de datos de alta calidad.

En la entrevista se explica que el valor de Scale consiste en que «los datos se venden cada vez más caros», porque las personas que emplea son cada vez mejores, las tareas son cada vez más complejas y los clientes necesitan cada vez más demostraciones y control de calidad de alto nivel. Para los modelos grandes y para los robots, los buenos datos ya no son solo etiquetado de bajo costo, sino el resultado conjunto de talento de primer nivel, procesos, relationship, sistemas de inspección de calidad y conocimiento del cliente. Si Meta quiere controlar la materia prima básica de la competencia futura en IA, buscará incorporar esta capacidad de datos a su propia estructura mediante una compra.

\begin{figure}[H]
\centering
\includegraphics[width=0.94\textwidth]{meta-scale-data.png}
\caption{La señal de la adquisición de Scale AI por parte de Meta: las empresas de diez billones del futuro necesitan dominar la capacidad de datos para IA. Diagrama conceptual propio, elaborado a partir de la conversación de 00:48:55--00:53:57.}
\end{figure}

\begin{knowledgebox}{Lectura de la figura: por qué la capacidad de datos pasó de ser subcontratada a ser estratégica}
En la figura se distinguen tres niveles de valor: el proceso de producción de datos, la demostración humana y la red de talento. Al principio, el etiquetado podía subcontratarse; cuando las tareas pasan a ser agentic, de robótica, de razonamiento complejo y de entornos de RL de alta calidad, la producción de datos se parece cada vez más a una capacidad central de investigación y desarrollo. La señal de Meta/Scale es esta: quien logre organizar a las mejores personas para definir los datos, filtrarlos y elevar su nivel de exigencia (bar) estará más cerca de la capacidad de los modelos de la siguiente etapa.
\end{knowledgebox}

\subsection{Los buenos datos son el producto de la diversidad de escenarios y de demostraciones de alta calidad}

Esta sección parte de una afirmación central: los buenos datos son el producto de dos cosas; la primera, escenarios diversos; la segunda, demostraciones humanas de alta calidad. Para los modelos grandes, la diversidad puede consistir en escenarios de tareas como matemáticas, programación, derecho o medicina; la demostración humana puede ser los pasos para resolver un problema, los criterios de evaluación o preguntas de examen de alta calidad. Para los robots, la diversidad está en hogares, hoteles, fábricas, restaurantes, anaqueles, cocinas, manijas de puertas y distintos objetos; la demostración humana, en cambio, consiste en trayectorias de movimiento, diseño de tareas, juicio sobre los fallos y diseño de la reward.

Además, el ponente sostiene que el mejor maestro no solo demuestra, sino que plantea problemas. Esta analogía explica por qué el trabajo con datos es cada vez más caro: no se trata de pedirle a alguien que haga una demostración cualquiera, sino de que diseñe problemas capaces de inducir el progreso del modelo, proporcione feedback de alta calidad y convierta los fallos en tareas de entrenamiento para la siguiente ronda. En el caso de los robots, un buen equipo de datos debe entender los escenarios físicos y también el entrenamiento de modelos, y además debe ser capaz de operar procesos de colaboración entre personas y máquinas a gran escala.

\begin{importantbox}{Fórmula de la calidad de los datos}
En el contexto de este episodio, los datos de alta calidad pueden entenderse así:
\[
\mathrm{Data\ Quality}\approx \mathrm{Diversity\ of\ Scenarios}\times \mathrm{Quality\ of\ Human\ Guidance}
\]
La diversidad de escenarios determina cuánto del mundo ha visto el modelo, y la guía humana de alta calidad determina si esos mundos pueden convertirse en señales de aprendizaje efectivas. Si falta cualquiera de las dos, los datos se empobrecen.
\end{importantbox}

\subsection{Cuellos de botella actuales de los datos sintéticos}

Esta sección vuelve a los datos sintéticos propiamente dichos. Xie Chen reconoce que, desde los primeros principios, el Sim2Real gap siempre existirá, porque lo sintético y lo real son distintos por naturaleza. Lo que importa no es eliminar el gap, sino saber a partir de qué magnitud del gap los datos resultan útiles y cómo reducirlo de manera continua. Esto exige que una empresa de datos no se limite a entregar datos, sino que también pueda llevar los algoritmos a la práctica, hacer despliegues reales y cerrar el ciclo de evaluación; de lo contrario, nunca sabrá dónde están las carencias.

Los cuellos de botella de los datos sintéticos pueden dividirse en tres etapas. En la etapa Real2Sim hay que aumentar la diversidad y la precisión de la información obtenida del mundo real; en la etapa Simulation hay que mejorar el motor físico subyacente, los recursos, los escenarios y la eficiencia en paralelo; en la etapa Sim2Real hay que mejorar la generalización y la efectividad de las políticas entrenadas. Por eso Guanglun (光轮), al mismo tiempo que produce datos, necesita entender los algoritmos de robótica de sus clientes y su puesta en práctica real.

\begin{warningbox}{Una empresa de datos no puede limitarse a entregar archivos}
Si una empresa de datos sintéticos solo se encarga de generar datos y el cliente entrena el modelo por su cuenta, a la empresa de datos le resultará difícil saber dónde está el gap real. Para mejorar de forma continua la calidad de los datos, debe participar en la evaluación del modelo, el despliegue real y el análisis de fallos.
\end{warningbox}

\subsection{El dilema de Physical Intelligence}

Esta sección examina la valoración que se hace en la entrevista de Physical Intelligence. Xie Chen considera que PI va muy adelante en modelos fundacionales corporizados, pero enfrenta un dilema respecto a la simulación: su discurso externo enfatiza los datos reales, mientras que internamente es muy consciente de que la simulación es necesaria. Esta contradicción afecta la atracción de talento, porque a los mejores especialistas en simulación les cuesta unirse a una organización que públicamente no da importancia a la simulación.

El valor didáctico de esta discusión es el siguiente: la simulación no es un pequeño componente de investigación, sino la combinación de cuatro cosas: tecnología, ingeniería, algoritmos y operación. Los equipos de investigación suelen dominar los algoritmos y los artículos, pero la simulación a gran escala requiere además plataformas de ingeniería, procesos operativos, inspección de calidad humana y captura física. Precisamente porque es pesada, sucia y compleja, se convierte en una oportunidad independiente dentro de la cadena industrial.

\begin{knowledgebox}{Términos clave: las cuatro cosas de la capacidad de simulación}
\begin{tabularx}{\textwidth}{p{0.18\textwidth}p{0.36\textwidth}X}
\toprule
Capacidad & Contenido concreto & Por qué la investigación por sí sola no basta \\
\midrule
Tecnología & Motor físico, renderizado, recursos y escenarios & Requiere acumulación de base a largo plazo. \\
Ingeniería & Plataformas a gran escala, cadena de herramientas, API y eficiencia en paralelo & Tanto el entrenamiento como las entregas al cliente exigen sistemas estables. \\
Algoritmos & Real2Sim, Sim2Real, RL y evaluación de calidad & Hay que convertir los datos en beneficios para el modelo. \\
Operación & Captura con humanos en el ciclo, inspección de calidad, retroalimentación de clientes y estandarización & La información del mundo físico no puede obtenerse de forma totalmente automática. \\
\bottomrule
\end{tabularx}
\end{knowledgebox}

\subsection{Resumen de la sección}

La discusión sobre Meta/Scale extiende el EP109 de los robots a toda la industria de la IA: la competencia futura girará cada vez más en torno a la capacidad de producir datos. Esto vale en especial para la inteligencia corporizada, porque los datos de alta calidad requieren diversidad del mundo y, además, demostraciones humanas, parámetros físicos, plataformas de ingeniería y un ciclo cerrado de evaluación.
